% Abstract and Bill Synopsis. One paragraph, under 250 words, no citations
% (every citation is anchored in the body). Rendered inside the abstract box.

The VVUQ Physical AI Oncology Trial Bill (H.R. 9510) is a proposed United States
law requiring an automated verification, validation, and uncertainty
quantification (VVUQ) process to clear robot-patient interaction code before that
code is generated or executed in a Physical AI oncology clinical trial. Its
thesis is rigor over the code itself: verification ahead of generation is the
patient-safety and efficacy control, and repository-based AI models synthesise
VVUQ code generation and execution into public, reproducible documents that
accelerate legislation. The evidence runs in three strokes. First, two author
studies show large language models are appropriate for code verification: one
proves the gate, not generation, is the decision-bearing work, accepting one of
six candidates while blocking five and escalating one, with the verification
fraction held at exactly 1.0 across 51 of 51 tests; the other proves practical
VVUQ generation and execution against fourteen external standards and two
clinical baselines, with 172 of 172 tests passing, 64 in a ten-gate assurance
suite, a 1{,}790-line four-entrant comparison, and a 1{,}001-row non-repetitive
humanoid sensor stream, both regenerated byte for byte from seed 20260525.
Second, a national platform supplies the Unification and Physical AI Standard
Levels and five trial, sponsor, and mobile simulations that make the
robot-patient case concrete. Third, the bill binds these to current federal and
state law. Codifying verification before generation keeps the United States first
in surgical-humanoid VVUQ while making oncology trials faster, less expensive,
and safer. The Unitree H2-Surgical 1.0 is hypothetical; every supporting number
is a simulation result.
